\workshoptitle{Bring your own laptop}{macOS and Linux}

\begin{context}
Download the event kit, install uv and Python, and create a folder for your bot.
The ZIP includes the client kit, its dependencies and all handouts. You do not
need Git, a compiler or a copy of the Arena server.
\end{context}

\lessonstep{Step 1}{Download and extract the kit}
Open the \href{https://github.com/Eastbridge-Academy/rps-client-kit/releases/latest}{GitHub release page}
(address below) and choose \textbf{rps-event-kit.zip} under Assets. Extract it
and find the folder containing \textbf{start.sh}, \textbf{wheelhouse} and
\textbf{requirements.txt}. Keep these files together; your bot will live in a
separate folder.

{\small\url{https://github.com/Eastbridge-Academy/rps-client-kit/releases/latest}\par}

\lessonstep{Step 2}{Install uv and Python}
Open Terminal. If \code{uv --version} already works, skip the first command.
Otherwise, install uv:

\begin{lstlisting}[style=shell]
curl -LsSf https://astral.sh/uv/install.sh | sh
\end{lstlisting}

Close Terminal and open it again, then install Python for the workshop:

\begin{lstlisting}[style=shell]
uv python install 3.12
\end{lstlisting}

These downloads need internet access. After this step, installing the kit and
playing local matches use the downloaded files. Submitting your bot and viewing
the tournament still need a connection.

\lessonstep{Step 3}{Create your bot project}
In Terminal, type \code{cd} followed by a space, drag the extracted folder
containing \textbf{start.sh} into the window, and press Enter. On Linux, you can
also right-click the folder and choose Open in Terminal. Then run:

\begin{lstlisting}[style=shell]
sh ./start.sh "$HOME/rps-bot"
cd "$HOME/rps-bot"
source .venv/bin/activate
rps-cli test
rps-cli validate
\end{lstlisting}

Open \textbf{bot.py} in your editor. Select the Python interpreter in
\textbf{rps-bot/.venv} if your editor asks. Each new terminal needs the
\code{cd} and \code{source} commands again.

\begin{checkpoint}
\textbf{Join the event:} return to \emph{Submit your bot} on page 1.
The facilitator supplies the event token. Run all bot commands from your
\textbf{rps-bot} folder.
\end{checkpoint}
